\newcounter{dois}
\stepcounter{dois}
\mysubsec{Exercício \thedois}
\bigline
\textbf{Sejam $(X, \mcal F, \mu)$ um espaço mensurável e $A \in \mcal F$. 
Mostre que a aplicação $\lambda: \mcal F \rightarrow \xbR$ dada por 
$\lambda(E)\coloneq \mu(E \cap A)$ é uma medida.}
\begin{proof}[\textbf{Dem:}]
Ora, sabendo que $\mu$ é uma medida, temos 
$\lambda(\emptyset) = \mu(\emptyset \cap A) = \mu(\emptyset) = 0$. Mais ainda, 
visto que para todo $B \in \mathcal F$, $\mu(B) \ge 0$, temos que dado $E \in \mathcal F$,
$E \cap A \in \mcal F$, ou seja, 
$$\lambda(E) = \mu(E \cap A) \ge 0.$$ 
Para verificar a propriedade de $\sigma$-aditividade, considere ${E_n}$ uma família em 
$\mathcal F$, dois a dois disjuntos. Como consequência da definição de ${E_n}$, temos que 
$\{E_n \cap A\}$ também é dois a dois disjuntos. Com isso,
\begin{align*}
   \lambda\left(\bigcup\limits_{n \in \mathbb N} E_n\right) &= 
   \mu\left(\left(\bigcup\limits_{n \in \mathbb N} E_n\right)\cap A\right) = 
   \mu\left(\bigcup\limits_{n \in \mathbb N} (E_n \cap A)\right),\\
   &=\sum_{n \in \mathbb N}\mu(E_n \cap A)=\sum_{n\in\mbb N}\lambda(E_n).
\end{align*}
Desse modo, segue que $\lambda$ é uma medida, como queríamos.
\end{proof}

\bigline
\newpage
\stepcounter{dois}
\mysubsec{Exercício \thedois}
\bigline
\textbf{Sejam $(X, \mcal F)$ um espaço mensurável $\mu_1, \cdots, \mu_n$ medidas em 
$(X, \mcal F)$ e $a_1, \cdots, a_n$ reais não negativos. Mostre que a aplicação 
$\lambda: \mcal F \rightarrow \xbR$ dada por 
\begin{align*}
  \lambda(E) = \sum_{k=1}^{n} a_k\mu_k(E)
\end{align*}
é uma medida em $(X, \mcal F)$.}
\begin{proof}[\textbf{Dem:}]
Visto que para todo $k \in \mbb N_{\le n}$, temos que $\mu_k$ é uma medida em $(X, \mathcal F)$,
então, para cada uma das medidas, segue que:
\begin{enumerate}
  \item $\mu_k(\emptyset) = 0$; 
  \item $\mu_k(E) \ge 0, \text{ para todo } E \in \mcal F$; 
  \item Vale a $\sigma$-aditividade para sequências de conjuntos disjuntos.
\end{enumerate}
Desse modo,
\begin{align*}
  \lambda(\emptyset) = \sum_{k=1}^{n} a_k\mu_k(\emptyset) \stackrel{(1)}{=} 
  \sum_{k=1}^{n} a_k \cdot0 = \sum_{k=1}^n 0 = 0. 
\end{align*}
Mais ainda, visto que $a_1, \dots, a_n$ é uma sequência de números não negativos,
  temos por $(2)$, que $a_k \mu_k \ge 0$ para todo $k \in \mbb N_{\le n}$. Ou seja 
\begin{align*}
  \lambda(E) = \sum_{k=1}^n a_k \mu_k(E) \ge 0, \text{ para todo } E \in \mcal F.
\end{align*}
  Por fim, dada uma sequência $\{E_j\}_{j \in \mbb N}$ em $\mcal F$ disjunto dois a dois,
  temos, 
  \begin{align*}
    \lambda\left(\bigcup\limits_{j\in \mbb N} E_j \right) &= 
    \sum_{k=1}^{n} \left[a_k \mu_k\left(\bigcup\limits_{j\in \mbb N} E_j\right)\right]
    \stackrel{(3)}{=}\sum_{k=1}^{n} \left[a_k \sum_{j \in \mbb N}\mu_k(E_j)\right]\\ 
    &= \sum_{k=1}^{n} \left[\sum_{j \in \mbb N}a_k\mu_k(E_j)\right]
    \end{align*}
Como a soma em $k$ é finita, podemos trocar a ordem das somas, obtendo
\begin{align*}
    &= \sum_{j \in \mbb N} \left[\sum_{k \in \mbb N_{\le n}}a_k\mu_k(E_j)\right]
    = \sum_{j \in \mbb N} \lambda(E_j).
  \end{align*}
  Com isso, temos que $\lambda$ é uma medida em $(X, \mcal F)$.
\end{proof}

\bigline
\newpage 
\stepcounter{dois}
\mysubsec{Exercício \thedois}
\bigline 
\textbf{
  Sejam $(X, \mcal F)$ um espaço mensurável e $(\mu_n)_{n \in \mbb N}$ uma sequência 
  de medidas em $(X, \mcal F)$ tal que $\mu_k(X) = 1$, para todo $k \in \mbb N$. 
  Mostre que a aplicação $\lambda: \mcal F \rightarrow \xbR$ dada por 
\begin{align*}
  \lambda(E) = \sum_{k \in \mbb N} \frac{\mu_k(E)}{2^k},
\end{align*}
é uma medida e $\lambda(X) = 1$.
}
\begin{proof}[\textbf{Dem:}] Vejamos que $\lambda(\emptyset) = 0$. 
Ora, $\lambda(\emptyset) = \sum_{k \in \mbb N} \frac{\mu_k(\emptyset)}{2^k}$. Já que 
  $\mu_k(\emptyset) = 0$ para todo $k \in \mbb N$, temos $\lambda(\emptyset) = 0$.
  Daqui, visto que para cada $E \in \mcal F$ e $k \in \mbb N$ temos $\mu_k(E)/2^k \ge 0$, 
então, para todo $E \in \mcal F$,
  \begin{align*}
    \lambda(E) = \sum_{k \in \mbb N} \frac{\mu_k(E)}{2^k} \ge 0.
  \end{align*}
  Com isso, vejamos que $\lambda$ é $\sigma$-aditiva. Dada uma sequência 
  $\{E_j\}_{j \in \mbb N}$ em $\mcal F$ disjunto dois a dois, temos, 
  \begin{align*}
    \lambda\left(\bigcup E_j\right) = 
    \sum_{k \in \mbb N} \frac{\mu_k \left(\bigcup E_j\right)}{2^k}.
  \end{align*}
  Já que vale a $\sigma$-aditividade para todo $\mu_k$, temos,
  \begin{align*}
    \lambda\left(\bigcup E_j\right) = 
    \sum_{k \in \mbb N}\left(\sum_{j \in \mbb N} \frac{\mu_k (E_j)}{2^k} \right).
  \end{align*}
  Como $\mu_k(E_j)/2^k \ge 0$ para cada $j,k \in \mbb N$,
 \begin{align*}
    \lambda\left(\bigcup E_j\right) = 
    \sum_{j \in \mbb N}\left(\sum_{k \in \mbb N} \frac{\mu_k (E_j)}{2^k} \right) = 
    \sum_{j \in \mbb N}\lambda(E_j).
  \end{align*}
 Por fim, visto que $\mu_k(X) = 1$ para todo $k \in \mbb N$, 
\begin{align*}
  \lambda(X) = \sum_{k \in \mbb N} \frac{\mu_k(X)}{2^k} 
  = \sum_{k \in \mbb N} \frac{1}{2^k} = 1. 
\end{align*}
Com isso, temos o que queríamos.
\end{proof}

\bigline
\newpage 
\stepcounter{dois}
\mysubsec{Exercício \thedois}
\bigline
\textbf{Seja $(\mbb N, \mcal P(\mbb N))$ e $(a_n)_{n \in \mbb N}$ uma sequência 
de números reais estendidos não negativos. Defina $\mu: \mcal P(\mbb N) \rightarrow \xbR$
por $\mu(\emptyset) = 0$ e, para $E \neq \emptyset$,
\begin{align*}
  \mu(E) = \sum_{n \in E} a_n.
\end{align*}
Mostre que $\mu$ é uma medida em $\mbb N$. Reciprocamente, mostre que toda medida 
em $(\mbb N, \mcal P(\mbb N))$ é obtida desta maneira, para alguma sequência 
$(a_n)$ em $\xbR$.
}
